% Generado con se-ramify. Para incluirlo en vuestra plantilla: \input{...}
% Necesita graphicx, longtable, booktabs, float, xcolor y fvextra.
% se-ramify: an image at its natural size, shrunk to fit the text block when bigger.
\providecommand{\seramifyfigura}[1]{%
  \sbox0{\includegraphics{#1}}%
  \ifdim\wd0>\linewidth
    \includegraphics[width=\linewidth,height=0.78\textheight,keepaspectratio]{#1}%
  \else\ifdim\ht0>0.78\textheight
    \includegraphics[height=0.78\textheight,keepaspectratio]{#1}%
  \else\usebox0\fi\fi}

\section*{Ejemplo — BiblioYa — Entrega 2 — UML, stakeholders y visión}
\addcontentsline{toc}{section}{Ejemplo — BiblioYa — Entrega 2 — UML, stakeholders y visión}

\emph{Exportado con se-ramify el 23/10/2026.}

\subsection*{Contenido de la entrega}

{\small
\begin{longtable}{@{}>{\raggedright\arraybackslash}p{0.453\linewidth}>{\raggedright\arraybackslash}p{0.061\linewidth}>{\raggedright\arraybackslash}p{0.409\linewidth}@{}}
\toprule
\textbf{Apartado} & \textbf{Peso} & \textbf{Diagramas} \\
\midrule
\endhead
Todos los diagramas UML: casos de uso con especificaciones, clases y secuencia & 60 \% & E2 · Casos de uso; E2 · Diagrama de clases de diseño; E2 · Secuencia — Reservar libro \\
Mapa de stakeholders (uno de los dos formatos, mínimo seis) & 20 \% & E2 · Mapa de stakeholders \\
Visión del proyecto según el modelo del "elevator pitch" & 20 \% & E2 · Visión — elevator pitch \\
\bottomrule
\end{longtable}}

\subsection*{1. E2 · Casos de uso}

\emph{Diagrama de sistema / casos de uso}

\begin{figure}[H]
\centering
\seramifyfigura{01-e2-casos-de-uso.png}
\caption{E2 · Casos de uso}
\end{figure}

\subsubsection*{CU-01 Buscar libro}

{\small
\begin{longtable}{@{}>{\raggedright\arraybackslash}p{0.233\linewidth}>{\raggedright\arraybackslash}p{0.716\linewidth}@{}}
\toprule
\textbf{Campo} & \textbf{Contenido} \\
\midrule
\endhead
Código & CU-01 \\
Nombre & Buscar libro \\
Actor principal & Socio \\
Actores secundarios &  \\
Descripción & El socio busca libros del catálogo por título, autor o ISBN y ve si hay ejemplares disponibles. \\
Disparador & El socio escribe en el buscador. \\
Precondiciones & Ninguna (la búsqueda es pública). \\
Flujo principal & 1. El socio escribe un título, autor o ISBN.\newline{}2. El sistema muestra los libros que coinciden, con el número de ejemplares disponibles.\newline{}3. El socio abre la ficha de un libro.\newline{}4. El sistema muestra sus ejemplares y su estado. \\
Flujos alternativos & 2a. No hay resultados: el sistema lo indica y sugiere revisar la ortografía. \\
Postcondiciones & Ninguna (no cambia datos). \\
Prioridad & Alta \\
Frecuencia & Cientos de veces al día \\
\bottomrule
\end{longtable}}

\subsubsection*{CU-02 Reservar libro}

{\small
\begin{longtable}{@{}>{\raggedright\arraybackslash}p{0.232\linewidth}>{\raggedright\arraybackslash}p{0.717\linewidth}@{}}
\toprule
\textbf{Campo} & \textbf{Contenido} \\
\midrule
\endhead
Código & CU-02 \\
Nombre & Reservar libro \\
Actor principal & Socio \\
Actores secundarios &  \\
Descripción & El socio reserva un libro que ahora está prestado para recogerlo cuando se devuelva. \\
Disparador & El socio pulsa «Reservar» en la ficha de un libro. \\
Precondiciones & El socio ha iniciado sesión y no tiene sanciones. \\
Flujo principal & 1. El socio busca el libro (CU-01).\newline{}2. El sistema muestra la ficha con sus ejemplares.\newline{}3. El socio pulsa «Reservar».\newline{}4. El sistema registra la reserva y muestra la posición en la cola. \\
Flujos alternativos & 3a. Hay un ejemplar disponible: el sistema ofrece reservarlo para recoger hoy.\newline{}4a. El socio ya tiene 3 reservas: el sistema lo impide y lo avisa. \\
Postcondiciones & La reserva queda registrada y el socio recibirá un aviso al devolverse el libro. \\
Prioridad & Media \\
Frecuencia & Unas 20 veces al día \\
\bottomrule
\end{longtable}}

\subsubsection*{CU-03 Prestar libro}

{\small
\begin{longtable}{@{}>{\raggedright\arraybackslash}p{0.232\linewidth}>{\raggedright\arraybackslash}p{0.717\linewidth}@{}}
\toprule
\textbf{Campo} & \textbf{Contenido} \\
\midrule
\endhead
Código & CU-03 \\
Nombre & Prestar libro \\
Actor principal & Bibliotecario \\
Actores secundarios & Socio \\
Descripción & El bibliotecario presta un ejemplar a un socio en el mostrador. \\
Disparador & Un socio entrega un libro en el mostrador para llevárselo. \\
Precondiciones & El bibliotecario ha iniciado sesión. \\
Flujo principal & 1. El bibliotecario identifica al socio (CU-04).\newline{}2. El bibliotecario lee el código del ejemplar.\newline{}3. El sistema comprueba que el socio puede pedir y el ejemplar está disponible.\newline{}4. El sistema registra el préstamo con fecha de devolución a 15 días.\newline{}5. El sistema imprime o envía el comprobante. \\
Flujos alternativos & 3a. El socio tiene una sanción o ya tiene 5 préstamos: el sistema lo impide y muestra el motivo.\newline{}3b. El ejemplar está reservado por otro socio: el sistema lo indica y no lo presta. \\
Postcondiciones & El ejemplar queda prestado al socio con su fecha de devolución. \\
Prioridad & Alta \\
Frecuencia & Unas 150 veces al día \\
\bottomrule
\end{longtable}}

\subsubsection*{CU-04 Identificar socio}

{\small
\begin{longtable}{@{}>{\raggedright\arraybackslash}p{0.232\linewidth}>{\raggedright\arraybackslash}p{0.717\linewidth}@{}}
\toprule
\textbf{Campo} & \textbf{Contenido} \\
\midrule
\endhead
Código & CU-04 \\
Nombre & Identificar socio \\
Actor principal & Bibliotecario \\
Actores secundarios &  \\
Descripción & El sistema reconoce al socio por su carné para operar en su nombre. \\
Disparador & El bibliotecario pasa el carné por el lector. \\
Precondiciones & El socio tiene carné en vigor. \\
Flujo principal & 1. El bibliotecario pasa el carné por el lector.\newline{}2. El sistema muestra los datos del socio, sus préstamos y sus sanciones. \\
Flujos alternativos & 1a. El carné no se lee: el bibliotecario escribe el número de socio.\newline{}2a. El carné está caducado: el sistema lo indica. \\
Postcondiciones & El socio queda identificado en el mostrador. \\
Prioridad & Alta \\
Frecuencia & Cada préstamo y devolución \\
\bottomrule
\end{longtable}}

\subsubsection*{CU-05 Devolver libro}

{\small
\begin{longtable}{@{}>{\raggedright\arraybackslash}p{0.232\linewidth}>{\raggedright\arraybackslash}p{0.717\linewidth}@{}}
\toprule
\textbf{Campo} & \textbf{Contenido} \\
\midrule
\endhead
Código & CU-05 \\
Nombre & Devolver libro \\
Actor principal & Bibliotecario \\
Actores secundarios & Socio \\
Descripción & El bibliotecario registra la devolución de un ejemplar y, si hay reservas, avisa al siguiente socio. \\
Disparador & Un socio devuelve un libro en el mostrador. \\
Precondiciones & El ejemplar está prestado. \\
Flujo principal & 1. El bibliotecario lee el código del ejemplar.\newline{}2. El sistema cierra el préstamo con la fecha de hoy.\newline{}3. El sistema comprueba si hay retraso.\newline{}4. El sistema marca el ejemplar como disponible. \\
Flujos alternativos & 3a. Hay retraso: el sistema calcula la sanción (1 día sin préstamo por día de retraso) y la muestra.\newline{}4a. Hay una reserva: el sistema lo marca como reservado y avisa al socio (CU-06). \\
Postcondiciones & El préstamo queda cerrado y el ejemplar disponible o reservado. \\
Prioridad & Alta \\
Frecuencia & Unas 150 veces al día \\
\bottomrule
\end{longtable}}

\subsubsection*{CU-06 Notificar disponibilidad}

{\small
\begin{longtable}{@{}>{\raggedright\arraybackslash}p{0.232\linewidth}>{\raggedright\arraybackslash}p{0.717\linewidth}@{}}
\toprule
\textbf{Campo} & \textbf{Contenido} \\
\midrule
\endhead
Código & CU-06 \\
Nombre & Notificar disponibilidad \\
Actor principal & Servicio de correo \\
Actores secundarios &  \\
Descripción & El sistema avisa al socio de que el libro que reservó ya está disponible. \\
Disparador & Se devuelve un ejemplar con reservas (CU-05). \\
Precondiciones & Hay una reserva y el libro se acaba de devolver. \\
Flujo principal & 1. El sistema toma la reserva más antigua.\newline{}2. El sistema envía un correo al socio con el plazo para recogerlo (3 días). \\
Flujos alternativos & 2a. El correo falla: el sistema lo reintenta en una hora y lo anota. \\
Postcondiciones & El socio está avisado y el ejemplar apartado para él. \\
Prioridad & Media \\
Frecuencia & Unas 20 veces al día \\
\bottomrule
\end{longtable}}

\subsection*{2. E2 · Diagrama de clases de diseño}

\emph{Diagrama de clases}

\begin{figure}[H]
\centering
\seramifyfigura{02-e2-diagrama-de-clases-de-diseno.png}
\caption{E2 · Diagrama de clases de diseño}
\end{figure}

\subsection*{3. E2 · Secuencia — Reservar libro}

\emph{Diagrama de secuencia}

\begin{figure}[H]
\centering
\seramifyfigura{03-e2-secuencia-reservar-libro.png}
\caption{E2 · Secuencia — Reservar libro}
\end{figure}

\subsection*{4. E2 · Mapa de stakeholders}

\emph{Mapa de stakeholders}

\begin{figure}[H]
\centering
\seramifyfigura{04-e2-mapa-de-stakeholders.png}
\caption{E2 · Mapa de stakeholders}
\end{figure}

\textbf{Producto / proyecto:} App de préstamos

{\scriptsize\setlength{\tabcolsep}{3pt}
\begin{longtable}{@{}>{\raggedright\arraybackslash}p{0.114\linewidth}>{\raggedright\arraybackslash}p{0.129\linewidth}>{\raggedright\arraybackslash}p{0.139\linewidth}>{\raggedright\arraybackslash}p{0.067\linewidth}>{\raggedright\arraybackslash}p{0.089\linewidth}>{\raggedright\arraybackslash}p{0.096\linewidth}>{\raggedright\arraybackslash}p{0.094\linewidth}>{\raggedright\arraybackslash}p{0.063\linewidth}>{\raggedright\arraybackslash}p{0.095\linewidth}@{}}
\toprule
\textbf{Stakeholder} & \textbf{Rol} & \textbf{Poder/Influencia} & \textbf{Interés} & \textbf{Estrategia} & \textbf{Anillo} & \textbf{Sector} & \textbf{Actitud} & \textbf{Notas} \\
\midrule
\endhead
Dirección de la biblioteca & Cliente: encarga y paga el proyecto & 90 \% & 85 \% & Gestionar de cerca & Es vital involucrarlo & Decisores & A favor & Aprueba el alcance y el presupuesto. \\
Vicerrectorado & Financia la biblioteca & 82 \% & 30 \% & Mantener satisfechos & Es necesario involucrarlos & Decisores & Neutral & Solo interviene si cambia el presupuesto. \\
Bibliotecarios & Usuarios internos diarios & 45 \% & 90 \% & Mantener informados & Es vital involucrarlo & Clientes & En contra & Temen más trabajo durante la migración. \\
Estudiantes & Usuarios finales & 30 \% & 80 \% & Mantener informados & Es necesario involucrarlos & Clientes & A favor &  \\
Profesores & Recomiendan bibliografía & 35 \% & 55 \% & Mantener informados & Estaría bien involucrarlos & Clientes & Neutral &  \\
Servicio de informática & Mantiene servidores y la red & 70 \% & 60 \% & Gestionar de cerca & Es necesario involucrarlos & Proveedores & Neutral & Necesitamos su visto bueno para integrar el catálogo. \\
Proveedor del catálogo & Licencia del software actual (OPAC) & 60 \% & 25 \% & Mantener satisfechos & Estaría bien involucrarlos & Proveedores & En contra & Puede cobrar por abrir su API. \\
Delegado de protección de datos & RGPD & 75 \% & 30 \% & Mantener satisfechos & Estaría bien involucrarlos & Influyentes & Neutral &  \\
Consejo de estudiantes & Representa a los alumnos & 25 \% & 45 \% & Monitorizar & Informar por cortesía & Influyentes & A favor &  \\
Otras bibliotecas de la red & Préstamo interbibliotecario & 15 \% & 20 \% & Monitorizar & Informar por cortesía & Proveedores & Neutral &  \\
\bottomrule
\end{longtable}}

\subsubsection*{Por estrategia (matriz poder/influencia – interés)}

\begin{itemize}
\item \textbf{Gestionar de cerca} (Champions · Drivers): Dirección de la biblioteca, Servicio de informática
\item \textbf{Mantener satisfechos} (Context setters · Blockers): Vicerrectorado, Proveedor del catálogo, Delegado de protección de datos
\item \textbf{Mantener informados} (Defenders · Subjects): Bibliotecarios, Estudiantes, Profesores
\item \textbf{Monitorizar} (Crowd · Bystanders): Consejo de estudiantes, Otras bibliotecas de la red
\end{itemize}

\subsection*{5. E2 · Visión — elevator pitch}

\emph{Visión — elevator pitch}

\begin{figure}[H]
\centering
\seramifyfigura{05-e2-vision-elevator-pitch.png}
\caption{E2 · Visión — elevator pitch}
\end{figure}

\begin{quote}
Para los socios de las bibliotecas municipales que quieren saber si un libro está disponible y llevárselo sin colas, BiblioYa es un servicio web y móvil de préstamos que permite buscar, reservar y renovar préstamos en dos clics. A diferencia del catálogo en línea actual, que solo informa, nuestro producto avisa cuando el libro vuelve, gestiona la cola de reservas y recuerda las devoluciones.
\end{quote}

{\small
\begin{longtable}{@{}>{\raggedright\arraybackslash}p{0.421\linewidth}>{\raggedright\arraybackslash}p{0.528\linewidth}@{}}
\toprule
\textbf{Plantilla} & \textbf{Contenido} \\
\midrule
\endhead
Para (cliente objetivo) & los socios de las bibliotecas municipales \\
Que (necesidad u oportunidad) & quieren saber si un libro está disponible y llevárselo sin colas \\
Producto (nombre) & BiblioYa \\
Es (categoría, con su artículo) & un servicio web y móvil de préstamos \\
Que (beneficio clave, razón para comprarlo) & permite buscar, reservar y renovar préstamos en dos clics \\
A diferencia de (alternativa / competidor) & el catálogo en línea actual, que solo informa \\
Nuestro producto (diferenciador) & avisa cuando el libro vuelve, gestiona la cola de reservas y recuerda las devoluciones \\
\bottomrule
\end{longtable}}
